\section{Further research questions and concrete next targets}\label{sec:research}

The following questions concern exact identities and structural reductions.
They are not assertions that numerical non-detection proves an obstruction.

\paragraph{1. Classify summation families admitting complete jet closure.}
For a general gamma-ratio sequence of ${}_{p+1}F_p$ type, the same
Bernoulli subtraction produces a normally convergent remainder, but the
unsubtracted endpoint need not have a gamma-product evaluation.
Determine precisely which classical summation submanifolds (for example,
balanced or well-poised parameter constraints) have a meromorphic gamma
product whose complete resonant jets yield finite Stieltjes closure.
A useful first target is a fixed Dougall-type summation with one excess
parameter and a fully explicit residue polynomial.  The essential issue is
an exact summation law, not merely existence of asymptotic coefficients.

\paragraph{2. Establish a multivariate resonance calculus without singular quotients.}
Formula \eqref{eq:EN} handles all zeros of $\rho_N$ for the present
one-parameter spectral problem.  For several simultaneous denominator
resonances, develop a Taylor algebra based on entire reciprocal-gamma
factors and finite head corrections.  Determine whether the resulting
constant terms depend on the path of a multivariate limit, and state the
normalization before introducing a multivariate finite part.  This would
prevent the loss of terms illustrated by \eqref{eq:piwarning}.

\paragraph{3. Identify the exact special-value algebra of rational harmonic-tail families.}
Theorem~\ref{thm:tail} gives a continuous $x$-parameter extension of a
multiple-zeta-star evaluation.  At rational $x$, decompose
$\zeta(2k+1,x)$ into cyclotomic or character coordinates, and express the
left-hand gamma weights and tails in a compatible alphabet.
The task is to obtain short explicit identities at denominators
$3,4,5,6$, not to claim minimal depth or independence.
The half-parameter formulas \eqref{eq:tail3}--\eqref{eq:tail7}
provide a model with no residual multiple constant on the right.

\paragraph{4. Construct mixed-order tail identities.}
The balanced variation $a=x+t,b=x-t$ generates only even tail power sums
$\zeta(2r,n+x)$.  Independent variations of $a,b$ generate both parities
and mixed Bell polynomials.  Determine which prescribed combinations of
odd and even tails have a single-zeta evaluation, which have finite
products of zeta values, and which retain nontrivial hypergeometric data.
The exact balanced identity \eqref{eq:balanced} provides a two-parameter
starting point, with all coefficient interchanges already justified.

\paragraph{5. Compress the primitive algebra at higher resonances.}
Theorem~\ref{thm:jets} permits nonpositive-integer zeta derivatives up to
$1-N$.  At small $(N,m)$ these can sometimes be rewritten in log-gamma,
Barnes-type, or normalized negative-polygamma coordinates.  Produce
canonical normalizations and change-of-basis formulas through, say,
$N=4$, $m=3$, including their rational-argument distribution laws.
A representation in a smaller \emph{specified} algebra is a meaningful
exact target; a claim that the remaining constants are genuinely new
requires an additional independence theorem.

\paragraph{6. Compare the fixed-shift and local-$z$ counterterm bases exactly.}
The present subtraction uses $(n+c)^{-1-u-j}$ and hence Lerch functions.
Classical continuation near $z=1$ instead uses powers of $1-z$ and
logarithms.  Determine finite connection matrices between these two
counterterm systems at arbitrary resonance and jet order, keeping every
constant term.  A complete comparison would sharpen the priority map
with the classical partial-sum literature and could explain additional
simplifications of \eqref{eq:alljets}.

\paragraph{7. Separate this calculus from the outstanding Gaussian certificates.}
The canonical $S_6$ and revised $S_8$ conjectures still require exact
relations in their stated colored-polylogarithm coordinates.  A useful
next attempt is to introduce a gamma-ratio parameter into an integral
representation of those coordinates and test whether the derivatives
land in a finite family controlled by a proved summation identity.
Any success must produce an explicit analytic identity or a replayable
exact relation certificate.  The present endpoint results alone do not
supply that missing certificate.

\paragraph{8. Formalize the finite coefficient layer and its analytic interface.}
The recursion, reference transport, and coefficient differential law are
polynomial identities suitable for exact formal verification.  The
analytic interface is a compact-normal convergence theorem for the
subtracted series.  Formalizing both layers would turn the present
regression tests into a machine-checked dependency chain, but no such
formalization is included or claimed here.

\section{Conclusion}
A single classical gamma-ratio summation, paired with a fixed Hurwitz
subtraction basis, gives a surprisingly extensive exact calculus.
Its zeroth-order resonances collapse to one digamma expression; its
higher spectral coefficients introduce Stieltjes constants in a controlled
finite way; balanced numerator variation gives single-zeta evaluations
of entire harmonic-tail families.  Reference differentiation telescopes,
while the generating variable supplies polylogarithmic endpoint
subtractions and explicit normalized primitives.

The important distinctions are structural.  An ordinary convergent
bracketed sum is not an unnormalized finite part.  A zero reciprocal-gamma
value is not a zero jet.  A finite closure formula is not an independence
result.  With these distinctions retained, the identities are directly
usable as a new, proved bridge within the ProveIt special-function
programme.
